% 原书第 8.8 节。
\section{大语言模型运维}
\label{sec:ch08-llmops}
随着人们对大语言模型的兴趣迅速增长，大语言模型运维（large language model operations，LLMOps）成为机器学习运维（MLOps）的一个自然分支。把大语言模型应用部署到生产环境时，许多问题与一般机器学习模型的部署问题相似。这里主要介绍与大语言模型有关的内容。有过其他模型运维经验的读者会熟悉其中不少方法。

当一种新思想迅速传播时，MLOps 和 LLMOps 往往成为流行词，具体涵盖什么也众说纷纭。这里采用较广义的理解，介绍能将机器学习能力部署到生产环境的多种工具与流程，包括数据和提示管理、迭代训练与工作流编排（\refcite{ch08-ref-064}{Oladele, 2024}）。多数方法旨在提高效率或降低风险，也有方法兼顾两者。这关系到能否从新的机器学习能力中获得价值。许多人曾把大量时间耗在模型上，给终端用户带来的好处很少，有时甚至有负面影响。

\subsection{LLMOps 工具与方法}
MLOps 及其后的 LLMOps 受到广泛重视，相应的解决方案市场迅速增长，出现了许多工具和产品。以下各节概述 LLMOps 生态，说明构建和部署大语言模型应用时会遇到哪些难题，以及常见的解决方式。图~\ref{fig:ch08-llmops}概述工具类别及其与生产应用的交互关系。

\sourcefigure{fig08_07.pdf}{一个基础聊天机器人及其 LLMOps 架构。提示模板经反复改进，并纳入版本控制。大语言模型也采用为该应用专门训练的适配器，因此需要跟踪实验并将训练好的组件发布到生产环境。提示模板和训练组件均通过 CI/CD 框架进入已部署的应用，生产指标则上报给监控服务。可以利用反馈进一步改进提示和训练数据。}{fig:ch08-llmops}

图的上部从提示工程、评估接到源码管理，从数据标注、训练流水线、实验跟踪接到模型注册库。下部依次展示用户界面、应用后端与推理服务器。后端用系统提示模板组织请求，可查询外部数据或服务，或直接响应用户。服务器负责请求路由、批处理，以及运行带适配器的预训练模型。示例中的用户查询订单，提供订单号 NX1234 后，收到商品缺货、预计于 2024 年 9 月 18 日发货的答复。

\subsubsection{实验跟踪}
把模型训练到生产部署所需的质量，需要反复迭代。特别是在开发新能力的早期，通常要进行大量试验。整理各轮结果并跟踪哪些模型表现最好，很快会变得繁重（Jakub \refcite{ch08-ref-032}{Czakon, 2024}）。日后可能还要找回其他细节，例如所用的具体数据集或所需的训练时间。许多实验跟踪工具可以满足这些需求。通常只需在训练流水线的代码中加入几条简单命令，相关信息便会自动记录。此类工具一般有功能完善的用户界面，提供多种指标和可视化方式，以便保存和分析实验。多人参与项目、希望查看彼此的工作时，这一点对团队协作特别有用。

模型注册库常与实验跟踪工具配合使用。模型注册库是保存通过实验跟踪工具创建的各个模型的仓库，也可直接上传未经实验跟踪的模型。模型注册库通常允许为模型添加自定义标签。下游流程可以据此自动识别哪些模型已可部署，或触发其他工作流。

\subsubsection{版本控制}
源码管理是软件项目中广泛采用的做法，同样适用于大语言模型。使用 Git 等代码仓库管理大语言模型的训练与评估代码，方式大体与管理其他代码库相同。以下还有一些仅靠代码仓库难以解决的版本管理需求。

\begin{itemize}
\item \textbf{模型版本管理。}构建大语言模型及其应用时，要多次迭代训练与调优。必须知道投入生产的是哪一版模型，并能找到生成该版本时用的确切代码和数据。否则，模型表现异常时很难调试并找出原因。采用实验跟踪工具和模型注册库后，这些工作大多已有工具支持。即便开发流程还没有那么成熟，无法使用先进工具，也必须以某种形式管理模型版本。
\item \textbf{数据版本管理。}训练数据往往是项目中变化最大的部分。常常要花大量时间确定最合适的数据类型，而用于编写代码的时间少得多。如果没有数据版本记录，模型就无法回滚至先前状态，相当于丢掉了以往的工作历史。
\item \textbf{提示版本管理。}提示和提示模板也是大语言模型系统的重要部分，可能在应用生命周期中发生很大变化。提示经常与应用代码一起存储，但这种做法未必始终合适。提示模板的行为通常类似概率模型，其验证方法因而与测试确定性代码的方式大不相同。单独管理提示可以简化开发，使团队能快速迭代提示，而无需每次都发布并部署较大组件的新版本。
\end{itemize}

\subsubsection{部署}
成熟软件开发生命周期中，许多常用于持续集成和持续交付或部署（CI/CD）的工具也适合部署大语言模型功能。这些流程旨在自动构建和测试待发布的新组件，并落实以下做法。

\begin{itemize}
\item 每次发布新版本都必须运行并通过同一组测试，以减少功能退化的可能。
\item 在通常模拟生产环境的集中环境中验证所有组件，统一各开发人员所用的验证条件。
\item 构建和发布过程应完全自动化且可重复，避免人工误操作使部署版本与通过测试的版本不同。
\item 组件投入生产后若表现未达要求，应能够迅速回滚至先前版本。
\end{itemize}

这种系统可以自动评估模型和报告指标，从而阻止表现欠佳的模型被误部署到生产环境。

\subsubsection{监控}
有许多工具可以监控已部署模型的表现，通常采用实时评估并汇总相关指标的方式。如果模型一开始就不符合预期，或其表现随时间发生变化，系统可以发出告警。例如，输出总体上比过去更短或更长，可能说明用户行为已经改变，也可能有其他原因。此时值得调查是否要相应调整模型或提示模板等组件。监控工具除了观察模型表现，还可防范模型幻觉和提示注入攻击等问题（\refcite{ch08-ref-004}{Ama, 2023}）。

可观测性与监控密切相关，两个术语有重叠，常被互换使用，市面上多数 LLMOps 方案也同时提供这两类功能。粗略区分的话，监控根据系统的汇总表现找出问题或有待改善的方面，可观测性则提供机制来进一步调查。例如，监控工具可能提示某天大语言模型的响应比平时慢。缺乏适当的可观测性时，可能难以找出原因。有了相应工具，就可以找出个别输入，逐步跟踪它在系统中的处理过程，定位瓶颈或故障。

\subsection{缩短迭代周期}
前面的章节讨论过人类反馈对于提高模型质量的重要性。人类反馈对基础模型的对齐和指令调优有价值，用户反馈对应用同样有价值。团队越快收集生产系统的结果，并用这些结果构建和发布更新，就越能为客户提供价值。改进这项工作流程，还可以迅速应对无法预见的问题。

另一个应缩短模型部署周期的理由是，用户反馈常常无可替代（\refcite{ch08-ref-012}{Burgess, 2021}）。数据科学家和开发人员做大量测试，也许能修正许多模型缺陷。不过，在处理真实用户输入之前，无法保证会得到预期结果。很多团队曾在新技术上投入数月乃至数年的努力，最终却在推出产品时受挫。避免这种局面的一种方式，是尽快推出最早可用的模型，同时准备好基础设施，以便取得新数据后可以迅速更新或回滚。ChatGPT 发布不久就出现了这种情况。谷歌和微软等主要竞争对手感到有必要尽快推出自己的聊天功能。好奇的用户很快开始攻击这些模型，从奇怪的响应中找乐趣，也暴露了模型的一些问题（\refcite{ch08-ref-065}{Orf, 2023}）。这些公司迅速处理早期问题，产品口碑最终没有受到太大影响。

\subsubsection{自动重新训练}
在某些情况下，可以直接用生产环境输入中的样例更新模型。例如，许多平台允许用户标记喜欢或不喜欢的内容。这类反馈可直接并入标注数据集，供下一轮训练使用。如果生产模型已较成熟，获取新数据、运行训练流水线、验证结果和部署新版本都可以全自动进行。需要快速适应新趋势的应用值得争取做到这一点，但实现起来并不容易。控制措施不足时，训练不充分的模型可能进入生产环境。应尽早考虑相应风险及降低这些风险所需的工作，再决定是否采用完全自动化的训练。达到这一成熟程度往往需要长期投入建设 LLMOps 能力。

\subsubsection{人在回路中的流程}
自动重新训练不可行或没有必要时，还有其他提高效率的方法。数据需要由人类专家审查或标注时，可采用多种标注工具。一些工具还可按需运用主动学习或半监督技术加快标注。前面讨论过，从人类反馈中进行强化学习（RLHF）显著提升了大语言模型的表现，所需的昂贵人力投入因此有了回报。收集或生成针对模型弱点的高质量训练样例，影响很大。样例的质量通常随人工投入增加。组织若无法为标注工作配备足够人员，也可将其外包，许多公司专门提供这类服务。

较新的数据或质量更可靠的数据往往会带来改进，可以迅速得到适于生产的模型，也有例外。模型架构可能是限制因素，这就需要进一步探索和研究，解决应用的弱点。实验跟踪和版本管理在这里很有帮助，特别是多人从不同角度研究同一个问题时。人工工作越多，越容易发生结果放错位置、数据集被改动等错误。还可以建立一个有代表性的黄金数据集供评估使用，并在整个迭代周期中保持固定，以便比较不同实验。

\subsection{风险管理}
用大语言模型或任何机器学习模型辅助完成传统上由人承担的任务，本身就涉及许多风险。设想一家公司被发现利用人工智能，以歧视性理由拒绝向某些人提供住房或就业机会，这可能造成严重影响。在这类环境中部署大语言模型时，需要特别认真地评估并尽量减少风险。

\begin{keypoint}
\textbf{风险评估与缓解}

根据大语言模型应用的功能进行充分的风险评估，是投入生产的重要一步。应考虑以下风险的影响。
\begin{itemize}
\item \textbf{侵犯隐私。}在正常使用过程中，应用是否会处理可能被意外展示给用户的受保护数据？
\item \textbf{安全漏洞。}应用设计和功能是否有可能导致数据丢失的漏洞？除了传统应用安全，还要考虑大语言模型所特有的提示注入攻击等风险。
\item \textbf{偏见和有害性。}应了解模型生成带有偏见或有害输出的风险。有些应用领域有明确的法律或监管规定，在这些领域尤其需要谨慎。
\end{itemize}
完成全面的风险评估后，可能需要采取缓解措施，包括以下几项。
\begin{itemize}
\item \textbf{人在回路中的监督。}许多请求量小、敏感度高或风险较大的使用场景，可能需要让人类专家参与。应用用户应能够在采取进一步行动之前核查大语言模型的输出。
\item \textbf{内容过滤。}应用可能需要实施内容过滤策略。可通过对齐调优或生成后过滤实现，例如采用关键词、模式匹配或分类方法，只将适宜的输出提供给用户。
\item \textbf{访问控制。}传统应用开发常有访问控制。在应用中使用大语言模型，可能还需要考虑其他角色权限。例如，对提示模板的操作权限可能只应授予担任提示工程师的开发人员。
\item \textbf{持续监控。}生产应用应记录正常的用户行为与交互情况，并持续监控。对大语言模型应用而言，了解用量、输出和表现有助于尽早发现问题。
\end{itemize}
\end{keypoint}

\subsubsection{模型治理}
前面已经看到，大语言模型仍面临若干问题，偏见只是其中之一。组织出于合理考虑，在受法律或监管约束的场景中使用机器学习时通常格外谨慎，尤其是医疗、金融和法律的多数领域。开发这类应用的技术人员应事先考虑要满足哪些要求，才能让利益相关者相信大语言模型的收益超过风险。许多组织制定了标准，要求对生产模型充分验证并记录。模型可解释性也可能至关重要。达到这些标准的工作量容易被低估。如果项目日程没有为此留出时间，就可能发生延误并增加成本。

模型卡（model card）是较常用的模型治理手段（\refcite{ch08-ref-059}{Mitchell et al., 2019}）。这项标准鼓励模型开发人员公开相关信息，降低模型被用于预期之外用途的风险。信息应以便于任何模型使用者查阅的方式呈现。不同模型适用的条目不尽相同。理想情况下，模型卡应说明推荐用途、已知限制和训练数据中的潜在偏见，也可以提供训练过程和不同基准上的评估指标。模型卡通常简明，技术细节较少。

\subsubsection{数据治理}
大语言模型快速兴起的一个重要原因是从网上抓取了大量数据。数据集迅速扩大后，整理或过滤特定类别的数据越来越难。大语言模型进入主流后，越来越多的人意识到自己的数据被用于未同意、甚至从未想过的用途，由此产生了矛盾。开发人员必须认真对待这些顾虑，以保护所在组织免受法律纠纷。

从社交媒体或论坛等网站提取数据时，可以先查阅相关条款。有些数据的使用规定可能不明确，也可能眼下似乎可以使用，日后却受到审查。因此，宜记录数据来源，包括每个数据集来自哪里、哪些模型使用过这些来源的数据。如果某些数据的使用因隐私、版权或其他问题而受到质疑，就可以减少损失，从存储中清除有关数据集，之后用排除这些数据的训练集训练模型。
